\begin{thebibliography}{99}
\raggedright

\bibitem{baker}
Alan Baker, \emph{Linear forms in the logarithms of algebraic numbers},
Mathematika \textbf{13} (1966), no.~2, 204--216.
\href{https://doi.org/10.1112/S0025579300003971}{doi:10.1112/S0025579300003971}.

\bibitem{green2026}
Ben Green, \emph{On Landau--Siegel zeros},
Hexagon preprint 2610.00311, version 2, 18 October 2026.
\url{https://hexagonmath.org/pdf/2610.00311v2}.

\bibitem{dlmf}
NIST Digital Library of Mathematical Functions, §§18.19 and 18.23,
\url{https://dlmf.nist.gov/18.19} and \url{https://dlmf.nist.gov/18.23}.

\bibitem{forresterwarnaar}
P. J. Forrester and S. O. Warnaar, \emph{The importance of the Selberg integral},
Bull. Amer. Math. Soc. (N.S.) \textbf{45} (2008), 489--534.
\href{https://doi.org/10.1090/S0273-0979-08-01221-4}{doi:10.1090/S0273-0979-08-01221-4}.

\bibitem{fulton}
William Fulton, \emph{Young Tableaux}, London Mathematical Society Student Texts,
vol.~35, Cambridge University Press, 1997.
\href{https://doi.org/10.1017/CBO9780511626241}{doi:10.1017/CBO9780511626241}.

\bibitem{heegner}
Kurt Heegner, \emph{Diophantische Analysis und Modulfunktionen},
Math. Z. \textbf{56} (1952), 227--253.
\href{https://doi.org/10.1007/BF01174749}{doi:10.1007/BF01174749}.

\bibitem{rosserschoenfeld}
J. Barkley Rosser and Lowell Schoenfeld,
\emph{Approximate formulas for some functions of prime numbers},
Illinois J. Math. \textbf{6} (1962), 64--94.
\href{https://doi.org/10.1215/ijm/1028998657}{doi:10.1215/ijm/1028998657}.

\bibitem{stark}
Harold M. Stark, \emph{A complete determination of the complex quadratic fields of class-number one},
Michigan Math. J. \textbf{14} (1967), 1--27.
\href{https://doi.org/10.1307/mmj/1028999653}{doi:10.1307/mmj/1028999653}.

\bibitem{taosketch}
Terence Tao, \href{https://terrytao.wordpress.com/2026/10/19/a-sketch-of-the-proof-of-non-existence-of-siegel-zeroes/}{\emph{A sketch of the proof of non-existence of Siegel zeroes}},
\emph{What's New}, 19 October 2026.

\bibitem{tao2026}
Terence Tao, \href{https://terrytao.wordpress.com/2026/10/20/an-experimental-ai-math-challenge-using-tenth-discriminant-bounds-to-explore-arguments-ruling-out-siegel-zeroes/}{\emph{An experimental AI math challenge: using tenth discriminant bounds to explore arguments ruling out Siegel zeroes}},
\emph{What's New}, 20 October 2026.

\end{thebibliography}
